% Slide 21 — Limitações atuais
\begin{frame}{Limitações atuais}
  \begin{itemize}
    \item Avaliação concentrada no regime \highlight{\textit{Lumpy}}
      (PB e BA); demais quadrantes de Syntetos-Boylan ainda não cobertos.
    \item Perfis \textit{Unstable Trend} ($n=18$) e \textit{High
      Vol.\ Seasonal} ($n=11$) têm $n<20$ -- evidência exploratória, não conclusiva.
    \item O oráculo por série é \highlight{limite teórico exploratório},
      não política implementável.
    \item O PSE ainda não foi treinado e validado em escala nacional.
    \item Horizonte fixo de 38 ciclos; parâmetros de custo e \textit{lead
      time} fixos no recorte avaliado.
    \item Previsão de demanda é sinal auxiliar -- não validada como critério
      de seleção isolado.
  \end{itemize}
  \note{
    É importante ser direto sobre o que ainda não está estabelecido. A
    avaliação atual está concentrada no regime Lumpy, o mais complexo
    operacionalmente, mas isso significa que a generalização para os
    regimes Smooth, Erratic e Intermittent ainda depende da expansão
    nacional. Dois dos três perfis operacionais analisados têm menos de 20
    séries e devem ser lidos como evidência exploratória. O oráculo por
    série é um limite teórico, não algo que se possa operacionalizar. E o
    PSE, apesar de formulado e testado nos rótulos disponíveis, ainda não
    foi treinado e validado na escala que permitiria afirmar capacidade de
    generalização. Essas limitações não invalidam os resultados -- elas
    delimitam exatamente o que pode ser afirmado nesta etapa.
  }
\end{frame}
